\begin{resumo}
Este trabalho avalia a capacidade de cinco Grandes Modelos de Linguagem (LLMs) — ChatGPT, Gemini, Claude, Meta AI e DeepSeek — em converter programas C sequenciais em implementações paralelas utilizando OpenMP. 
A avaliação foi conduzida sobre seis programas do benchmark PAMPAR, abrangendo desde reduções com iterações independentes até algoritmos com dependências de dados complexas. 
Cada modelo recebeu o mesmo prompt padronizado e o código-fonte sequencial, em uma única tentativa, sem refinamento iterativo. 
Os programas gerados foram avaliados quanto à compilação, à correção da saída e ao speedup obtido em uma máquina de memória compartilhada, utilizando as implementações de referência do PAMPAR como alvo de desempenho. 
Todos os programas gerados compilaram e produziram saída idêntica à versão sequencial. 
Em desempenho, os modelos de fronteira aproximaram-se da implementação de referência escrita por especialistas, enquanto os demais apresentaram perdas de eficiência concentradas nos programas com padrões de paralelização mais complexos. 
Os resultados indicam que LLMs atuais são capazes de gerar código paralelo funcionalmente correto para uma ampla variedade de programas em uma única tentativa, mas diferem de forma significativa em eficiência à medida que a complexidade da paralelização aumenta.

\vspace{\onelineskip}

\noindent
\textbf{Palavras-chave}: Grandes Modelos de Linguagem. Paralelização Automática. OpenMP. Computação de Alto Desempenho.
\end{resumo}
